\section{Golden Branch Orientation and Irrational Scanning}\label{sec:golden-branch}

\subsection{Irrational Orientation and Quasiperiodicity (H4.1)}

If the orientation of $E_{\parallel}$ relative to $\mathcal{L}$ is ``rational,'' then $\pi_{\parallel}(\mathcal{L})$ contains nontrivial periodic lattice structure, and $\Lambda(W,u)$ degenerates to a periodic family. To obtain nonperiodic but long-range ordered point sets, adopt the following replaceable assumption:
\begin{itemize}
  \item \textbf{(H4.1)} $(E_{\parallel},E_{\perp})$ is in irrational orientation relative to $\mathcal{L}$, such that the internal projection $\pi_{\perp}(\mathcal{L})$ is dense in an appropriate quotient space, and the model set belongs to the quasiperiodic class.
\end{itemize}
Under this condition, the patch statistics of $\Lambda(W,u)$ can be controlled by hull-space ergodicity in typical cases (Appendix \ref{app:general-position}).

\subsection{Scanning Fiber and Protocol Time (D4.2)}

To describe ``phenomena varying with time'' without introducing object-layer time coordinates, introduce a discrete scanning of internal phase:
\[
u_t:=u_0+t\beta \ \ (\mathrm{mod}\ \Gamma),\qquad t\in\ZZ_{\ge 0},
\]
where $\Gamma\subset E_{\perp}$ is the internal equivalence lattice (e.g., induced by translational symmetries of $\pi_{\perp}(\mathcal{L})$), and $\beta\in E_{\perp}$ is the scanning direction. Protocol time is the iteration index $t$.

\subsection{Equidistribution and Recurrence Scale (T4.3, H4.3)}

\begin{itemize}
  \item \textbf{(H4.3)} Assume that $\beta$ satisfies irrationality on $E_{\perp}/\Gamma$, such that $(u_t)$ is equidistributed.
\end{itemize}
Under this assumption, any integrable function $f$ satisfies time-averaging equal to space-averaging (ergodicity semantics in Appendix \ref{app:general-position}), so ``long-range statistics are reproducible'' obtains formal assurance.

\subsection{Anti-Resonance Condition of Golden Branch (D4.4, T4.4)}

In finite-depth scanning, near-resonance (spurious short-period artifacts from rational approximation) is controlled by the Diophantine properties of $\beta$. This paper defines:
\begin{itemize}
  \item \textbf{(D4.4)} A scanning direction (or its one-dimensional factor rotation number) is said to belong to the \textbf{golden branch} if its continued-fraction partial quotients achieve the minimum under complexity budget, and are thus the most weakly approximable within the same budget.
  \item \textbf{(T4.4)} In one-dimensional circle rotation, the golden-ratio-related rotation numbers achieve the uniformly optimal ``anti-resonance extremum'' in the sense of Hurwitz approximation constants; hence they delay the shortest recurrence time scale within a fixed complexity budget.
\end{itemize}
For the higher-dimensional case, verifiable checks can be obtained through one-dimensional factorization (Section \ref{sec:symbolic}) and multidimensional discrepancy estimation (Appendix \ref{app:golden-branch}).
